\documentclass[12pt,a4paper]{article}

% ─── PACKAGES ─────────────────────────────────────────────
\usepackage[utf8]{inputenc}
\usepackage[T1]{fontenc}
\usepackage[french]{babel}
\usepackage{amsmath,amssymb,amsthm}
\usepackage{mathtools}
\usepackage{booktabs}
\usepackage{array}
\usepackage{longtable}
\usepackage{graphicx}
\usepackage{xcolor}
\usepackage{hyperref}
\usepackage{enumitem}
\usepackage{geometry}
\usepackage{fancyhdr}
\usepackage{titlesec}
\usepackage{tcolorbox}
\usepackage{caption}
\usepackage{subcaption}
\usepackage{algorithm}
\usepackage{algpseudocode}
\usepackage{tikz}
\usetikzlibrary{arrows.meta, positioning, shapes.geometric, calc}

\geometry{margin=2.5cm}

% ─── COULEURS ─────────────────────────────────────────────
\definecolor{steelblue}{HTML}{4682B4}
\definecolor{darkblue}{HTML}{2C3E50}
\definecolor{lightgray}{HTML}{F5F5F5}
\definecolor{accentgreen}{HTML}{27AE60}
\definecolor{accentorange}{HTML}{E67E22}
\definecolor{accentred}{HTML}{E74C3C}

% ─── STYLE ────────────────────────────────────────────────
\hypersetup{
    colorlinks=true,
    linkcolor=steelblue,
    citecolor=steelblue,
    urlcolor=steelblue,
}

\titleformat{\section}
    {\Large\bfseries\color{darkblue}}{\thesection}{1em}{}
\titleformat{\subsection}
    {\large\bfseries\color{steelblue}}{\thesubsection}{1em}{}
\titleformat{\subsubsection}
    {\normalsize\bfseries\color{steelblue!80}}{\thesubsubsection}{1em}{}

\pagestyle{fancy}
\fancyhf{}
\fancyhead[L]{\small\color{gray}IFRS 9 Risk Cockpit --- Analyse CRO}
\fancyhead[R]{\small\color{gray}\thepage}
\renewcommand{\headrulewidth}{0.4pt}

% ─── ENCADRES ─────────────────────────────────────────────
\newtcolorbox{definitionbox}[1][]{
    colback=steelblue!5,
    colframe=steelblue,
    title=#1,
    fonttitle=\bfseries,
    rounded corners,
    boxrule=0.8pt,
}

\newtcolorbox{formulebox}[1][]{
    colback=lightgray,
    colframe=darkblue,
    title=#1,
    fonttitle=\bfseries,
    rounded corners,
    boxrule=0.8pt,
}

\newtcolorbox{interpretbox}[1][]{
    colback=accentgreen!5,
    colframe=accentgreen!70!black,
    title=#1,
    fonttitle=\bfseries,
    rounded corners,
    boxrule=0.8pt,
}

\newtcolorbox{warningbox}[1][]{
    colback=accentorange!5,
    colframe=accentorange,
    title=#1,
    fonttitle=\bfseries,
    rounded corners,
    boxrule=0.8pt,
}

% ─── DOCUMENT ─────────────────────────────────────────────
\begin{document}

\begin{titlepage}
\centering
\vspace*{3cm}
{\Huge\bfseries\color{darkblue} Analyse CRO --- 5 Couches\par}
\vspace{0.8cm}
{\Large\color{steelblue} AI Analyst / Virtual CRO\par}
\vspace{1.5cm}
{\large\color{gray} IFRS 9 Risk Cockpit\par}
\vspace{2cm}
{\large Pipeline analytique 2 passes : exploration, d\'etection de r\'egime,\\ conditionnement et synth\`ese narrative\par}
\vspace{3cm}
{\normalsize
\begin{tabular}{ll}
\textbf{Module} & \texttt{ifrs9\_cockpit.ai\_analyst} \\
\textbf{Orchestrateur} & \texttt{orchestrator.py} --- classe \texttt{CROAnalyst} \\
\textbf{Couches} & 5 modules Python (layer1 \`a layer5) \\
\textbf{Sortie} & \texttt{AnalyticsState} (dataclass auto-document\'ee) \\
\end{tabular}
}
\vfill
{\small\color{gray} Document g\'en\'er\'e depuis le code source --- confidentiel}
\end{titlepage}

\tableofcontents
\newpage

% ═══════════════════════════════════════════════════════════
\section{Architecture 5 couches}
% ═══════════════════════════════════════════════════════════

\subsection{Vue d'ensemble}

L'AI Analyst impl\'emente un pipeline analytique \`a 5 couches orchestr\'e en 2 passes
par la classe \texttt{CROAnalyst} (fichier \texttt{orchestrator.py}). Chaque couche
produit une composante de l'\'etat analytique \texttt{AnalyticsState}, un dataclass
auto-document\'e consomm\'e par le dashboard Streamlit (onglet 8 : Analyse CRO).

\begin{definitionbox}[Pipeline 2 passes]
\begin{enumerate}[label=\textbf{Passe \arabic*},leftmargin=2.5cm]
    \item \textbf{Exploration} (r\'egime inconnu) --- les 5 couches s'ex\'ecutent sans
          conditionnement. La couche~4 d\'etecte le r\'egime macro actuel.
    \item \textbf{Conditionnement} --- les 5 couches se r\'e-ex\'ecutent avec le r\'egime
          d\'etect\'e en passe~1. L'attribution factorielle (couche~2) et les trajectoires
          (couche~5) sont ajust\'ees par des multiplicateurs et drifts sp\'ecifiques au
          r\'egime.
\end{enumerate}
Apr\`es la passe~2, l'orchestrateur g\'en\`ere les \textbf{recommandations explicables}
(FR31) et la \textbf{synth\`ese narrative} (FR32).
\end{definitionbox}

\subsection{Diagramme d'architecture}

\begin{figure}[htbp]
\centering
\begin{tikzpicture}[
    layer/.style={
        rectangle, rounded corners=6pt, minimum width=12cm, minimum height=1.4cm,
        draw=steelblue, thick, fill=steelblue!8, font=\small\bfseries,
        text=darkblue, align=center
    },
    sublabel/.style={
        font=\scriptsize\color{gray}, align=center
    },
    orch/.style={
        rectangle, rounded corners=8pt, minimum width=13.5cm, minimum height=1.2cm,
        draw=darkblue, very thick, fill=darkblue!10, font=\normalsize\bfseries,
        text=darkblue
    },
    synth/.style={
        rectangle, rounded corners=6pt, minimum width=12cm, minimum height=1.2cm,
        draw=accentgreen!70!black, thick, fill=accentgreen!8, font=\small\bfseries,
        text=accentgreen!70!black
    },
    data/.style={
        ellipse, draw=accentorange, thick, fill=accentorange!10,
        font=\scriptsize, text=accentorange!80!black, align=center,
        minimum width=2.8cm, minimum height=0.9cm
    },
    myarrow/.style={-{Stealth[length=5pt]}, thick, steelblue},
    feedback/.style={-{Stealth[length=5pt]}, thick, accentred, dashed},
]

% Orchestrateur
\node[orch] (orch) at (0, 8.5) {Orchestrateur CROAnalyst --- 2 Passes};

% Couches
\node[layer] (L1) at (0, 6.8) {Couche 1 --- Croisement \& Asym\'etries};
\node[sublabel, below=-0.1cm of L1] (L1s) {\texttt{layer1\_crossing.py} --- 5 secteurs $\times$ 2 canaux, contributions marginales};

\node[layer] (L2) at (0, 5.0) {Couche 2 --- Allocation Proportionnelle \& Attribution Factorielle};
\node[sublabel, below=-0.1cm of L2] (L2s) {\texttt{layer2\_euler.py} --- allocation proportionnelle 10 cellules, deltas adaptatifs (M6)};

\node[layer] (L3) at (0, 3.2) {Couche 3 --- Reverse Stress Testing};
\node[sublabel, below=-0.1cm of L3] (L3s) {\texttt{layer3\_rst.py} --- tipping points univari\'es, RST joint Mahalanobis};

\node[layer] (L4) at (0, 1.4) {Couche 4 --- D\'etection de R\'egime};
\node[sublabel, below=-0.1cm of L4] (L4s) {\texttt{layer4\_regime.py} --- 5 r\'egimes canoniques, softmax sur similarit\'e cosinus};

\node[layer] (L5) at (0, -0.4) {Couche 5 --- Trajectoires Prospectives};
\node[sublabel, below=-0.1cm of L5] (L5s) {\texttt{layer5\_prospective.py} --- Ornstein-Uhlenbeck, risk appetite, early warning};

% Synthese
\node[synth] (synth) at (0, -2.2) {Synth\`ese : Recommandations (FR31) + Narrative (FR32)};

% Fleches verticales
\draw[myarrow] (orch.south) -- (L1.north);
\draw[myarrow] (L1.south) -- ++(0,-0.2) -- (L2.north);
\draw[myarrow] (L2.south) -- ++(0,-0.2) -- (L3.north);
\draw[myarrow] (L3.south) -- ++(0,-0.2) -- (L4.north);
\draw[myarrow] (L4.south) -- ++(0,-0.2) -- (L5.north);
\draw[myarrow] (L5.south) -- ++(0,-0.2) -- (synth.north);

% Feedback passe 2
\draw[feedback] (L4.east) -- ++(2.2,0) node[midway, above, font=\scriptsize\color{accentred}] {r\'egime} |- (orch.east);

% Donnees d'entree
\node[data] (dc) at (-7.5, 6.8) {result\_credit\\(ECLCalculator)};
\node[data] (dp) at (-7.5, 5.0) {result\_pe\\(PECalculator)};
\node[data] (dm) at (-7.5, 3.2) {macro\_params\\(5 variables)};

\draw[myarrow, accentorange] (dc.east) -- (L1.west);
\draw[myarrow, accentorange] (dp.east) -- (L2.west);
\draw[myarrow, accentorange] (dm.east) -- (L3.west);

\end{tikzpicture}
\caption{Architecture du pipeline AI Analyst --- 5 couches, 2 passes}
\label{fig:architecture}
\end{figure}

\subsection{Donn\'ees d'entr\'ee}

L'orchestrateur re\c{c}oit trois entr\'ees principales :

\begin{itemize}
    \item \texttt{result\_credit} : \texttt{DataFrame} produit par \texttt{ECLCalculator.calculate()},
          contenant par ligne : \texttt{sector}, \texttt{ecl\_weighted}, \texttt{ead},
          \texttt{rwa\_credit}, \texttt{pd\_12m}, \texttt{stage}.
    \item \texttt{result\_pe} : \texttt{DataFrame} produit par \texttt{PECalculator.calculate()},
          contenant : \texttt{sector}, \texttt{expected\_loss\_pe}, \texttt{rwa\_pe},
          \texttt{nav}, \texttt{nav\_drawdown}, \texttt{risk\_category}.
    \item \texttt{macro\_params} : dictionnaire des 5 variables macro\-\'economiques
          (\texttt{unemployment\_rate}, \texttt{gdp\_growth}, \texttt{interest\_rate},
          \texttt{hpi\_growth}, \texttt{inflation\_rate}).
\end{itemize}

\subsection{Structure de sortie : AnalyticsState}

La sortie compl\`ete est port\'ee par le dataclass \texttt{AnalyticsState}
(fichier \texttt{types.py}), qui expose des champs typ\'es pour chaque couche :

\begin{table}[htbp]
\centering
\caption{Champs de \texttt{AnalyticsState} par couche}
\begin{tabular}{clp{7.5cm}}
\toprule
\textbf{Couche} & \textbf{Champ} & \textbf{Description} \\
\midrule
1 & \texttt{asymmetry\_matrix} & DataFrame 5 secteurs : taux de perte credit vs PE, asym\'etrie \\
1 & \texttt{marginal\_contributions} & 10 cellules (secteur $\times$ canal), contribution marginale \\
2 & \texttt{proportional\_contributions} & Allocation proportionnelle 10 cellules (\texttt{proportional\_share}) \\
2 & \texttt{factor\_attribution} & Attribution factorielle OAT, 5 vars $\times$ 2 canaux \\
3 & \texttt{tipping\_points} & Seuils de basculement univari\'es (5 variables) \\
3 & \texttt{rst\_result} & Sc\'enario de rupture RST (dict) \\
3 & \texttt{rst\_distance} & Distance de Mahalanobis en $\sigma$ \\
4 & \texttt{regime} & \texttt{RegimeClassification} (5 probabilit\'es + r\'egime dominant) \\
5 & \texttt{trajectories} & 12 projections (3 trajectoires $\times$ 4 horizons) \\
5 & \texttt{risk\_appetite\_matrix} & 10 cellules avec feux tricolores \\
5 & \texttt{early\_warning} & Indicateurs d'alerte pr\'ecoce composites (5 secteurs) \\
--- & \texttt{recommendations} & Liste de \texttt{Recommendation} (FR31) \\
--- & \texttt{narrative} & Synth\`ese narrative CRO (FR32) \\
\bottomrule
\end{tabular}
\label{tab:analytics_state}
\end{table}


\newpage
% ═══════════════════════════════════════════════════════════
\section{Couche 1 : Analyse descriptive et croisements}
% ═══════════════════════════════════════════════════════════

\subsection{Objectif}

La couche~1 (\texttt{layer1\_crossing.py}, fonction \texttt{analyze\_crossings})
scanne les \textbf{asym\'etries credit vs Private Equity} par secteur et calcule
les \textbf{contributions marginales} de chaque cellule (secteur $\times$ canal)
au risque total du portefeuille.

\subsection{Matrice d'asym\'etrie}

Pour chaque secteur $s \in \{1, \ldots, 5\}$, on calcule :

\begin{formulebox}[Taux de perte normalis\'e]
\begin{equation}
    \ell_s^{\text{Credit}} = \frac{\text{ECL}_s}{\text{EAD}_s}, \qquad
    \ell_s^{\text{PE}} = \frac{\text{EL\_PE}_s}{\text{NAV}_s}
    \label{eq:loss_rates}
\end{equation}
\begin{equation}
    \text{Asym\'etrie}_s = \ell_s^{\text{PE}} - \ell_s^{\text{Credit}}
    \label{eq:asymmetry}
\end{equation}
\end{formulebox}

Une asym\'etrie positive indique que le canal PE est plus risqu\'e que le canal credit
pour ce secteur. L'intensit\'e de stress par secteur est calcul\'ee comme la somme pond\'er\'ee
des \'ecarts macro par les sensibilit\'es sectorielles :

\begin{equation}
    I_s = \sum_{v=1}^{5} \left| X_v - X_v^{\text{base}} \right| \cdot \gamma_{s,v}^{\text{credit}}
    \label{eq:stress_intensity}
\end{equation}

o\`u $\gamma_{s,v}^{\text{credit}}$ est la sensibilit\'e du secteur $s$ \`a la variable $v$,
d\'efinie dans \texttt{SectorConfig}.

\subsection{Contributions marginales}

Les contributions marginales reposent sur une allocation proportionnelle simple :

\begin{equation}
    c_{s,k} = \frac{R_{s,k}}{R_{\text{total}}}, \quad
    R_{\text{total}} = \sum_s \text{ECL}_s + \sum_s \text{EL\_PE}_s
    \label{eq:marginal}
\end{equation}

o\`u $k \in \{\text{Credit}, \text{PE}\}$ et $R_{s,k}$ est le montant de risque
de la cellule $(s, k)$. La propri\'et\'e de full allocation est garantie :
$\sum_{s,k} c_{s,k} = 1$.

\subsection{Risk Appetite --- Cadre tricolore}

Le cadre d'app\'etit au risque s'applique par cellule via la fonction
\texttt{compute\_risk\_appetite} (couche~5). Les seuils sont d\'efinis dans
\texttt{RiskAppetiteConfig} :

\begin{table}[htbp]
\centering
\caption{Seuils de l'app\'etit au risque (cadre FSB 2013)}
\begin{tabular}{lccc}
\toprule
\textbf{M\'etrique} & \textcolor{accentgreen}{\textbf{Vert}} & \textcolor{accentorange}{\textbf{Ambre}} & \textcolor{accentred}{\textbf{Rouge}} \\
\midrule
ECL/EAD (Credit) & $< 2{,}0\%$ & $[2{,}0\% ; 4{,}0\%[$ & $\geq 4{,}0\%$ \\
NAV Drawdown (PE) & $< 10\%$ & $[10\% ; 25\%[$ & $\geq 25\%$ \\
RAROC & $> 4\%$ & $[2\% ; 4\%]$ & $< 2\%$ \\
HHI & $< 1\,500$ & $[1\,500 ; 2\,500[$ & $\geq 2\,500$ \\
\bottomrule
\end{tabular}
\label{tab:risk_appetite}
\end{table}

\subsection{Indicateurs d'alerte pr\'ecoce (Early Warning)}

La fonction \texttt{compute\_early\_warning} calcule un score composite par secteur
via une pond\'eration de type AHP simplifi\'ee :

\begin{formulebox}[Score d'alerte pr\'ecoce]
\begin{equation}
    \text{EW}_s = 0{,}25 \cdot \widetilde{\text{PD}}_s
    + 0{,}15 \cdot \widetilde{\text{S2}}_s
    + 0{,}15 \cdot \widetilde{\text{S3}}_s
    + 0{,}20 \cdot \widetilde{\text{DD}}_s
    + 0{,}10 \cdot \widetilde{\text{Dist}}_s
    + 0{,}15 \cdot \widetilde{M}
    \label{eq:early_warning}
\end{equation}
o\`u chaque composante est normalis\'ee par sa valeur de r\'ef\'erence maximale :
$\widetilde{\text{PD}}_s = \min(1, \text{PD}_s / 0{,}20)$, \quad
$\widetilde{\text{S2}}_s = \min(1, \%\text{Stage2} / 0{,}20)$, etc.
\end{formulebox}

Les poids (somme~$= 1$) sont calibr\'es par jugement d'expert :
PD moyenne (0,25), Stage~2 (0,15), Stage~3 (0,15), drawdown PE (0,20),
distress PE (0,10), stress macro global (0,15).
Le signal est vert si $\text{EW} < 0{,}25$, ambre si $< 0{,}50$, rouge sinon.


\newpage
% ═══════════════════════════════════════════════════════════
\section{Couche 2 : Allocation proportionnelle du risque}
% ═══════════════════════════════════════════════════════════

\subsection{Objectif}

La couche~2 (\texttt{layer2\_euler.py}, fonction \texttt{decompose\_proportional})
d\'ecompose le risque total en contributions par cellule (secteur $\times$ canal)
et attribue les variations de risque aux 5 variables macro\-\'economiques.

\begin{warningbox}[Allocation proportionnelle (ex-``Euler'')]
\textbf{Audit RJ} : l'ancienne terminologie <<~Euler~>> a \'et\'e remplac\'ee par
<<~allocation proportionnelle~>>. La m\'ethode n'utilise \textbf{pas} de d\'eriv\'ees
partielles (Euler-Tasche, 1999) mais une simple r\`egle de trois :
\[
    \text{contribution}_i = \frac{\text{ECL}_i}{\text{ECL}_{\text{total}}}
\]
qui garantit la propri\'et\'e de \textbf{full allocation}
($\sum_i \text{contribution}_i = 1$) par construction.
L'alias \texttt{decompose\_euler} reste disponible pour la compatibilit\'e ascendante.
\end{warningbox}

\subsection{Allocation proportionnelle (10 cellules)}

Pour les 5 secteurs et 2 canaux (Credit, PE), le risque total est d\'efini comme :

\begin{equation}
    R_{\text{total}} = \underbrace{\sum_{s=1}^{5} \text{ECL}_s}_{\text{Credit}}
    + \underbrace{\sum_{s=1}^{5} \text{EL\_PE}_s}_{\text{Private Equity}}
    \label{eq:total_risk}
\end{equation}

La contribution proportionnelle de chaque cellule est :

\begin{formulebox}[Contribution proportionnelle]
\begin{equation}
    w_{s}^{\text{Credit}} = \frac{\text{ECL}_s}{R_{\text{total}}}, \qquad
    w_{s}^{\text{PE}} = \frac{\text{EL\_PE}_s}{R_{\text{total}}}
    \label{eq:proportional}
\end{equation}
Propri\'et\'e : $\displaystyle\sum_{s=1}^{5} \left( w_s^{\text{Credit}} + w_s^{\text{PE}} \right) = 1$
\end{formulebox}

Le DataFrame retourn\'e contient pour chaque cellule :
\texttt{sector}, \texttt{canal}, \texttt{risk\_amount} (montant ECL ou EL en EUR),
\texttt{proportional\_share} (contribution normalis\'ee), et \texttt{rwa} (RWA associ\'e).

\subsection{Attribution factorielle OAT}

L'attribution factorielle mesure la sensibilit\'e du risque \`a chaque variable
macro\-\'economique par la m\'ethode One-At-a-Time (OAT) avec des
\textbf{deltas adaptatifs} (correction M6).

\subsubsection{Deltas adaptatifs (M6)}

Plut\^ot qu'un delta fixe de 1~point de pourcentage pour toutes les variables,
chaque delta est calibr\'e \`a 1\% de la fourchette de chaque variable :

\begin{table}[htbp]
\centering
\caption{Deltas adaptatifs par variable (1\% de la fourchette)}
\begin{tabular}{lccr}
\toprule
\textbf{Variable} & \textbf{Fourchette} & \textbf{Delta adaptatif} \\
\midrule
\texttt{unemployment\_rate} & $[4 ; 15]$ & 0,11 \\
\texttt{gdp\_growth} & $[-5 ; 5]$ & 0,10 \\
\texttt{interest\_rate} & $[0 ; 10]$ & 0,10 \\
\texttt{hpi\_growth} & $[-10 ; 10]$ & 0,20 \\
\texttt{inflation\_rate} & $[0 ; 8]$ & 0,08 \\
\bottomrule
\end{tabular}
\label{tab:adaptive_deltas}
\end{table}

\subsubsection{Formule d'attribution}

Pour chaque variable $v$ et canal $k$, l'attribution est :

\begin{formulebox}[Attribution factorielle normalis\'ee]
\begin{equation}
    A_{v}^{k} = \sum_{s=1}^{5}
    \frac{|\Delta_v|}{\delta_v^{\text{adapt}}}
    \cdot \gamma_{s,v}^{k}
    \cdot \frac{R_{s}^{k}}{R_{\text{total}}^{k}}
    \label{eq:attribution}
\end{equation}
o\`u :
\begin{itemize}[nosep]
    \item $\Delta_v = X_v^{\text{current}} - X_v^{\text{base}}$ est l'\'ecart au baseline,
    \item $\delta_v^{\text{adapt}}$ est le delta adaptatif (Table~\ref{tab:adaptive_deltas}),
    \item $\gamma_{s,v}^{k}$ est la sensibilit\'e du secteur $s$ \`a la variable $v$ pour le canal $k$.
\end{itemize}
\end{formulebox}

\subsubsection{Ajustement r\'egime (passe 2)}

En passe~2, l'attribution factorielle est multipli\'ee par un facteur r\'egime
$m_{v,r}$ qui amplifie les variables dominantes du r\'egime d\'etect\'e $r$ :

\begin{equation}
    A_{v}^{k,\text{ajust}} = A_{v}^{k} \cdot m_{v,r}
    \label{eq:regime_mult}
\end{equation}

\begin{table}[htbp]
\centering
\caption{Multiplicateurs de r\'egime (extraits de \texttt{\_regime\_multiplier})}
\begin{tabular}{lccccc}
\toprule
\textbf{R\'egime} & \texttt{unemp} & \texttt{gdp} & \texttt{rate} & \texttt{hpi} & \texttt{infl} \\
\midrule
Crise financi\`ere & 1,5 & 1,3 & 1,0 & 1,4 & 1,0 \\
Stagflation & 1,0 & 1,3 & 1,2 & 1,0 & 1,5 \\
Rupture techno & 1,2 & 1,4 & 1,0 & 1,0 & 1,0 \\
Resserrement & 1,0 & 1,0 & 1,5 & 1,3 & 1,0 \\
Reprise & 0,8 & 0,8 & 1,0 & 1,0 & 1,0 \\
\bottomrule
\end{tabular}
\label{tab:regime_mult}
\end{table}

\begin{interpretbox}[Interpr\'etation]
En r\'egime de \textbf{crise financi\`ere}, le ch\^omage ($\times 1{,}5$),
les prix immobiliers ($\times 1{,}4$) et le PIB ($\times 1{,}3$) sont amplifi\'es,
refl\'etant la dynamique observ\'ee lors de la GFC 2008--09.
En \textbf{reprise}, les multiplicateurs sont $< 1$ (att\'enuation de l'attribution),
coh\'erent avec un risque en d\'ecroissance.
\end{interpretbox}


\newpage
% ═══════════════════════════════════════════════════════════
\section{Couche 3 : Reverse Stress Testing}
% ═══════════════════════════════════════════════════════════

\subsection{Objectif}

La couche~3 (\texttt{layer3\_rst.py}) remplit deux fonctions compl\'ementaires :

\begin{enumerate}
    \item \textbf{Seuils de basculement univari\'es} : pour chaque variable macro, identifier
          la valeur qui fait basculer l'ECL au-del\`a du seuil ambre.
    \item \textbf{Reverse Stress Test joint} : trouver le sc\'enario de rupture
          \textbf{minimal} (au sens de Mahalanobis) qui provoque un breach CET1.
\end{enumerate}

\subsection{Proxy ECL logit}

Le RST utilise un proxy rapide de l'ECL total (plutot que le recalcul complet
\texttt{ECLCalculator}) pour permettre le balayage param\'etrique en temps r\'eel :

\begin{formulebox}[Proxy ECL logit]
\begin{equation}
    \text{stress}_{\text{sum}} = \sum_{v=1}^{5}
    \frac{\Delta_v}{100} \cdot \bar{\gamma}_v \cdot
    \begin{cases}
        -1 & \text{si } v \in \{\texttt{gdp}, \texttt{hpi}\} \\
        +1 & \text{sinon}
    \end{cases}
    \label{eq:stress_sum}
\end{equation}
\begin{equation}
    \text{ECL}_{\text{proxy}} = \sigma\!\left(\text{logit}(r_{\text{base}})
    + \text{stress}_{\text{sum}} \cdot \Lambda\right) \cdot \text{EAD}_{\text{total}}
    \label{eq:ecl_proxy}
\end{equation}
o\`u $\sigma$ est la fonction logistique (expit), $r_{\text{base}} = \text{ECL} / \text{EAD}$,
$\Lambda = 6{,}0$ (\texttt{LOGIT\_AMPLITUDE}), et $\bar{\gamma}_v$ est la sensibilit\'e
moyenne pond\'er\'ee par la proportion sectorielle.
\end{formulebox}

Le retour est en \textbf{montant absolu EUR} (ratio $\times$ EAD), coh\'erent
avec le seuil de rupture du capital.

\subsection{Seuils de basculement univari\'es}

La fonction \texttt{find\_tipping\_points} balaye chaque variable dans sa direction
adverse avec un pas de 0,5~pp :

\begin{algorithm}[htbp]
\caption{Recherche de seuils de basculement univari\'es}
\label{alg:tipping}
\begin{algorithmic}[1]
\Require $f_{\text{ECL}}$ : proxy ECL, $\mathbf{X}$ : variables macro, $\tau$ : seuil ECL
\Ensure DataFrame de 5 tipping points
\For{chaque variable $v \in \{1, \ldots, 5\}$}
    \State $x_0 \gets X_v$ \Comment{valeur courante}
    \State $d \gets$ direction adverse de $v$ \Comment{+1 ou $-1$}
    \For{$x \gets x_0,\; x_0 + d \cdot 0{,}5,\; x_0 + d \cdot 1{,}0,\; \ldots$}
        \If{$x \notin [\text{min}_v, \text{max}_v]$}
            \State \textbf{break} \Comment{hors fourchette}
        \EndIf
        \State $\mathbf{X}' \gets \mathbf{X}$ avec $X_v' = x$
        \If{$f_{\text{ECL}}(\mathbf{X}') > \tau$}
            \State \Return tipping\_value $= x$, distance $= |x - x_0|$
        \EndIf
    \EndFor
    \State tipping\_value $= \varnothing$ (pas de basculement dans la fourchette)
\EndFor
\end{algorithmic}
\end{algorithm}

Le seuil $\tau$ par d\'efaut est le seuil ambre en montant absolu :
$\tau = \texttt{ecl\_ead\_amber} \times \text{EAD}_{\text{total}}$.
Le param\`etre \texttt{ecl\_threshold} (FR54) permet de le personnaliser.

\begin{table}[htbp]
\centering
\caption{Fourchettes de balayage par variable}
\begin{tabular}{lcccl}
\toprule
\textbf{Variable} & \textbf{Min} & \textbf{Max} & \textbf{Pas} & \textbf{Direction adverse} \\
\midrule
\texttt{unemployment\_rate} & 4,0 & 15,0 & 0,5 & $\uparrow$ \\
\texttt{gdp\_growth} & $-5{,}0$ & 5,0 & 0,5 & $\downarrow$ \\
\texttt{interest\_rate} & 0,0 & 10,0 & 0,5 & $\uparrow$ \\
\texttt{hpi\_growth} & $-10{,}0$ & 10,0 & 0,5 & $\downarrow$ \\
\texttt{inflation\_rate} & 0,0 & 8,0 & 0,5 & $\uparrow$ \\
\bottomrule
\end{tabular}
\label{tab:sweep_ranges}
\end{table}

\subsection{Reverse Stress Test joint}

La fonction \texttt{reverse\_stress\_test} cherche le facteur d'\'echelle
minimal $\alpha \in [1{,}0 ; 10{,}0]$ (pas de 0,1) tel que toutes les variables
sont stress\'ees simultan\'ement et l'ECL d\'epasse le seuil de rupture CET1.

\subsubsection{Construction du sc\'enario stress\'e}

Pour un facteur $\alpha$, les variables stress\'ees sont :

\begin{equation}
    X_v^{\text{stress}}(\alpha) = X_v^{\text{base}}
    + \frac{(X_v^{\max} - X_v^{\text{base}}) \cdot d_v \cdot (\alpha - 1)}{5}
    \label{eq:stress_params}
\end{equation}

o\`u $d_v \in \{+1, -1\}$ est la direction adverse de la variable $v$.

\subsubsection{Seuil de rupture}

Le seuil ECL de rupture est calibr\'e \`a 10\% du capital CET1 r\'eel :

\begin{equation}
    \text{ECL}_{\text{breach}} = \underbrace{\text{RWA}_{\text{credit}} \times \text{CET1}_{\text{target}}}_{\text{Capital CET1 r\'eel}} \times 0{,}10
    \label{eq:breach}
\end{equation}

Le param\`etre \texttt{target\_ecl} permet de personnaliser ce seuil (FR54).
Le capital est calcul\'e dynamiquement depuis le portefeuille r\'eel
(pas de constante statique).

\subsubsection{Distance de Mahalanobis (H6)}

\begin{definitionbox}[Distance de Mahalanobis]
La distance de Mahalanobis mesure la distance statistique entre le sc\'enario
de rupture et le baseline, en tenant compte de la structure de corr\'elation
des variables macro :
\begin{equation}
    d_M = \sqrt{\mathbf{x}^\top \boldsymbol{\Sigma}^{-1} \mathbf{x}}
    \label{eq:mahalanobis}
\end{equation}
o\`u $\mathbf{x} = \mathbf{X}^{\text{stress}} - \mathbf{X}^{\text{base}}$
est le vecteur d'\'ecart (5 dimensions) et $\boldsymbol{\Sigma}$ est la matrice
de covariance $5 \times 5$ construite \`a partir d'une corr\'elation de consensus expert
et de volatilit\'es historiques (cf.\ section~\ref{tab:scenarios}).
\end{definitionbox}

L'impl\'ementation dans \texttt{\_compute\_mahalanobis} utilise
\texttt{np.linalg.solve} au lieu de l'inversion explicite
$\boldsymbol{\Sigma}^{-1} \mathbf{x}$ pour une meilleure stabilit\'e num\'erique :

\begin{formulebox}[Impl\'ementation num\'erique stable]
Au lieu de calculer $\boldsymbol{\Sigma}^{-1} \mathbf{x}$ directement, on r\'esout
le syst\`eme lin\'eaire $\boldsymbol{\Sigma} \mathbf{y} = \mathbf{x}$ pour obtenir
$\mathbf{y} = \boldsymbol{\Sigma}^{-1} \mathbf{x}$, puis :
\begin{equation}
    d_M = \sqrt{\mathbf{x}^\top \mathbf{y}}
    \label{eq:solve}
\end{equation}
En cas de singularit\'e de $\boldsymbol{\Sigma}$ (exception \texttt{LinAlgError}),
un fallback diagonal (distance euclidienne normalis\'ee) est utilis\'e :
\begin{equation}
    d_M^{\text{fallback}} = \sqrt{\sum_{i=1}^{5} \frac{x_i^2}{\sigma_{ii}^2}}
    \label{eq:fallback}
\end{equation}
\end{formulebox}

La matrice $\boldsymbol{\Sigma}$ est construite dans \texttt{config.py} par la fonction
\texttt{\_build\_macro\_covariance()} \`a partir d'une matrice de corr\'elation de consensus
(\textit{Expert Judgment}) et de volatilit\'es historiques ($\Sigma_{ij} = \rho_{ij} \sigma_i \sigma_j$).
L'ordre des variables est fix\'e par \texttt{MACRO\_VARIABLES\_ORDER} :
(\texttt{unemployment\_rate}, \texttt{gdp\_growth}, \texttt{interest\_rate},
\texttt{hpi\_growth}, \texttt{inflation\_rate}).

\begin{interpretbox}[Interpr\'etation de la distance RST]
\begin{itemize}[nosep]
    \item $d_M < 2\sigma$ : rupture \`a proximit\'e --- risque \'elev\'e, confiance \textbf{low}
    \item $2\sigma \leq d_M < 5\sigma$ : rupture mod\'er\'ement \'eloign\'ee --- confiance \textbf{medium}
    \item $d_M \geq 5\sigma$ : rupture tr\`es \'eloign\'ee --- confiance \textbf{high}
\end{itemize}
Cette grille est utilis\'ee directement par l'orchestrateur pour g\'en\'erer le
niveau de confiance des recommandations.
\end{interpretbox}


\newpage
% ═══════════════════════════════════════════════════════════
\section{Couche 4 : D\'etection de r\'egime}
% ═══════════════════════════════════════════════════════════

\subsection{Objectif}

La couche~4 (\texttt{layer4\_regime.py}, fonction \texttt{classify\_regime})
classifie le r\'egime macro\-\'economique actuel avec des \textbf{probabilit\'es continues}
sur 5 r\'egimes canoniques, via un softmax appliqu\'e aux scores de similarit\'e.

\subsection{Les 5 r\'egimes canoniques}

\begin{definitionbox}[R\'egimes canoniques]
Cinq r\'egimes couvrent l'ensemble des configurations macro\-\'economiques pertinentes
pour le risque de cr\'edit :
\begin{enumerate}[nosep]
    \item \textbf{Crise financi\`ere} --- calibr\'e sur la GFC 2008--09
          (ch\^omage $+3$, PIB $-3$, taux $-1$, immobilier $-5$)
    \item \textbf{Stagflation} --- calibr\'e sur les ann\'ees 1970
          (ch\^omage $+2$, PIB $-2$, taux $+2$, inflation $+3$)
    \item \textbf{Rupture techno} --- sc\'enario hypoth\'etique
          (ch\^omage $+1$, PIB $-1{,}5$)
    \item \textbf{Resserrement} --- calibr\'e sur 2022--23
          (taux $+3$, immobilier $-3$, inflation $+1{,}5$)
    \item \textbf{Reprise} --- calibr\'e sur 2021
          (ch\^omage $-2$, PIB $+2$, immobilier $+2$)
\end{enumerate}
Source : BCE Statistical Data Warehouse, calibration p\'edagogique.
\end{definitionbox}

\subsection{Vecteur de stress}

Le vecteur de stress est calcul\'e comme l'\'ecart entre les variables macro
actuelles et le baseline, avec une \textbf{convention d'inversion} pour le PIB
et les prix immobiliers (une baisse correspond \`a un stress positif) :

\begin{equation}
    \mathbf{s} = \begin{pmatrix}
        X_{\text{unemp}} - X_{\text{unemp}}^{\text{base}} \\
        X_{\text{gdp}}^{\text{base}} - X_{\text{gdp}} \\
        X_{\text{rate}} - X_{\text{rate}}^{\text{base}} \\
        X_{\text{hpi}}^{\text{base}} - X_{\text{hpi}} \\
        X_{\text{infl}} - X_{\text{infl}}^{\text{base}}
    \end{pmatrix}
    \label{eq:stress_vector}
\end{equation}

\subsection{Classification par produit scalaire + softmax}

Le score de chaque r\'egime $r$ est le \textbf{produit scalaire} entre le vecteur
de stress actuel et la signature du r\'egime :

\begin{equation}
    z_r = \mathbf{s} \cdot \boldsymbol{\sigma}_r = \sum_{i=1}^{5} s_i \cdot \sigma_{r,i}
    \label{eq:dot_product}
\end{equation}

o\`u $\boldsymbol{\sigma}_r$ est le vecteur-signature du r\'egime $r$
(Table~\ref{tab:signatures}).

Les probabilit\'es sont obtenues par softmax avec temp\'erature $T = 2{,}0$
(\texttt{\_SOFTMAX\_SCALE}) :

\begin{formulebox}[Softmax avec temp\'erature]
\begin{equation}
    P(r) = \frac{\exp(T \cdot z_r)}{\displaystyle\sum_{r'=1}^{5} \exp(T \cdot z_{r'})}
    \label{eq:softmax}
\end{equation}
Le r\'egime d\'etect\'e est $r^* = \arg\max_r P(r)$.
Pour la stabilit\'e num\'erique, le calcul utilise le trick standard
$\exp(T z_r - \max_j(T z_j))$ avant normalisation.
\end{formulebox}

\begin{table}[htbp]
\centering
\caption{Signatures des 5 r\'egimes (espace stress)}
\label{tab:signatures}
\begin{tabular}{lccccc}
\toprule
\textbf{R\'egime} & \texttt{unemp} & \texttt{gdp} & \texttt{rate} & \texttt{hpi} & \texttt{infl} \\
\midrule
Crise financi\`ere & $+3{,}0$ & $+3{,}0$ & $-1{,}0$ & $+5{,}0$ & $0{,}0$ \\
Stagflation & $+2{,}0$ & $+2{,}0$ & $+2{,}0$ & $+1{,}0$ & $+3{,}0$ \\
Rupture techno & $+1{,}0$ & $+1{,}5$ & $+0{,}5$ & $0{,}0$ & $+0{,}5$ \\
Resserrement & $+0{,}5$ & $+0{,}5$ & $+3{,}0$ & $+3{,}0$ & $+1{,}5$ \\
Reprise & $-2{,}0$ & $-2{,}0$ & $0{,}0$ & $-2{,}0$ & $-0{,}5$ \\
\bottomrule
\end{tabular}
\end{table}

\begin{interpretbox}[Interpr\'etation]
\begin{itemize}[nosep]
    \item En l'absence de stress ($\mathbf{s} \approx \mathbf{0}$), les scores
          sont proches et les probabilit\'es quasi-uniformes ($\sim 20\%$ chacun).
    \item Sous stress de type \emph{Stagflation} (ch\^omage $+3$, PIB $-1$,
          inflation $+5$, taux $+5{,}5$), la signature Stagflation domine par
          produit scalaire et la probabilit\'e associ\'ee peut d\'epasser 80\%.
    \item La temp\'erature $T = 2{,}0$ offre un bon compromis entre discrimination
          (haute temp\'erature $\to$ winner-takes-all) et incertitude
          (basse temp\'erature $\to$ uniforme).
\end{itemize}
\end{interpretbox}


\newpage
% ═══════════════════════════════════════════════════════════
\section{Couche 5 : Trajectoires prospectives}
% ═══════════════════════════════════════════════════════════

\subsection{Objectif}

La couche~5 (\texttt{layer5\_prospective.py}) projette les variables macro
\`a T+3, T+6, T+9 et T+12 mois selon 3 sc\'enarios (favorable, central, adverse)
via un processus \textbf{d'Ornstein-Uhlenbeck} discret (mean-reversion).

\subsection{Processus d'Ornstein-Uhlenbeck (H10)}

\begin{definitionbox}[Processus O-U continu]
Le processus d'Ornstein-Uhlenbeck mod\'elise le retour \`a la moyenne d'une variable :
\begin{equation}
    \mathrm{d}X_t = \kappa\,(\theta - X_t)\,\mathrm{d}t + \sigma\,\mathrm{d}W_t
    \label{eq:ou_continuous}
\end{equation}
o\`u $\kappa > 0$ est la vitesse de retour \`a la moyenne, $\theta$ est l'\'equilibre
de long terme, $\sigma$ est la volatilit\'e annualis\'ee, et $W_t$ est un mouvement
brownien standard.
\end{definitionbox}

L'impl\'ementation utilise la discr\'etisation d'Euler avec un pas mensuel
$\Delta t = 1/12$ :

\begin{formulebox}[Discr\'etisation O-U (pas mensuel)]
\begin{equation}
    X_{t+1} = X_t + \kappa \cdot (\theta - X_t) \cdot \Delta t
    + \sigma \cdot \sqrt{\Delta t} \cdot \varepsilon
    \label{eq:ou_discrete}
\end{equation}
o\`u $\varepsilon \in \{-1{,}5 ;\; 0 ;\; +1{,}5\}$ est un choc \textbf{d\'eterministe}
d\'efinissant les 3 sc\'enarios :
\begin{itemize}[nosep]
    \item \textbf{Favorable} : $\varepsilon = -1{,}5$ (am\'elioration)
    \item \textbf{Central} : $\varepsilon = 0$ (neutre, retour \`a la moyenne pur)
    \item \textbf{Adverse} : $\varepsilon = +1{,}5$ (d\'et\'erioration)
\end{itemize}
\end{formulebox}

\subsection{Param\`etres O-U par variable}

Les param\`etres $(\kappa, \sigma, \theta)$ sont d\'efinis dans
\texttt{MACRO\_MEAN\_REVERSION} (fichier \texttt{config.py}).

\begin{warningbox}[Audit RJ -- d\'ecouplage $\theta$ / baseline]
$\theta$ repr\'esente l'\textbf{\'equilibre structurel de long terme}, ind\'ependant de la
conjoncture actuelle (\texttt{SCENARIO\_BASE}). Les valeurs sont calibr\'ees sur le consensus
OCDE/FMI long terme et les estimations structurelles (NAIRU, $r^*$ Laubach-Williams,
cible BCE). Le dictionnaire \texttt{MACRO\_STRUCTURAL\_EQUILIBRIUM} porte ces valeurs.
\end{warningbox}

\begin{table}[htbp]
\centering
\caption{Param\`etres Ornstein-Uhlenbeck par variable macro}
\begin{tabular}{lcccc}
\toprule
\textbf{Variable} & $\kappa$ (vitesse) & $\sigma$ (vol.) & $\theta$ (\'equilibre structurel) & Source \\
\midrule
\texttt{unemployment\_rate} & 0,5 & 0,8 & 6,5\% & NAIRU zone euro \\
\texttt{gdp\_growth} & 1,0 & 1,2 & 1,5\% & Potentiel UE \\
\texttt{interest\_rate} & 0,3 & 0,5 & 2,5\% & $r^*$ Laubach-Williams \\
\texttt{hpi\_growth} & 0,7 & 2,0 & 2,0\% & Inflation + productivit\'e \\
\texttt{inflation\_rate} & 0,8 & 0,6 & 2,0\% & Cible BCE \\
\bottomrule
\end{tabular}
\label{tab:ou_params}
\end{table}

\begin{interpretbox}[Lecture des param\`etres]
\begin{itemize}[nosep]
    \item Le \textbf{PIB} a la vitesse de retour la plus \'elev\'ee ($\kappa = 1{,}0$)
          et la volatilit\'e la plus forte ($\sigma = 1{,}2$), coh\'erent avec les cycles
          conjoncturels courts.
    \item Le \textbf{taux directeur} a la vitesse la plus faible ($\kappa = 0{,}3$),
          refl\'etant l'inertie de la politique mon\'etaire.
    \item L'\textbf{immobilier} est le plus volatil en niveau ($\sigma = 2{,}0$) avec
          une vitesse de retour mod\'er\'ee ($\kappa = 0{,}7$).
\end{itemize}
\end{interpretbox}

\subsection{Ajustement r\'egime en passe 2}

En passe~2, un \textbf{drift additionnel} d\'ependant du r\'egime d\'etect\'e est ajout\'e
aux projections :

\begin{equation}
    X_v^{\text{proj, ajust}}(T) = X_v^{\text{proj}}(T) + \delta_{v,r} \cdot \frac{T}{12}
    \label{eq:regime_drift}
\end{equation}

o\`u $\delta_{v,r}$ est le drift par an du r\'egime $r$ sur la variable $v$.

\begin{table}[htbp]
\centering
\caption{Drifts additionnels par r\'egime (par an)}
\begin{tabular}{lccccc}
\toprule
\textbf{R\'egime} & \texttt{unemp} & \texttt{gdp} & \texttt{rate} & \texttt{hpi} & \texttt{infl} \\
\midrule
Crise financi\`ere & $+0{,}5$ & $-0{,}5$ & --- & $-1{,}0$ & --- \\
Stagflation & --- & $-0{,}3$ & $+0{,}3$ & --- & $+0{,}5$ \\
Rupture techno & $+0{,}3$ & $-0{,}2$ & --- & --- & --- \\
Resserrement & --- & --- & $+0{,}5$ & $-0{,}5$ & --- \\
Reprise & $-0{,}3$ & $+0{,}3$ & --- & --- & --- \\
\bottomrule
\end{tabular}
\label{tab:regime_drifts}
\end{table}

\subsection{Format de sortie}

La fonction retourne un DataFrame de 12 lignes :
3 trajectoires $\times$ 4 horizons (T+3, T+6, T+9, T+12),
avec les colonnes \texttt{trajectory}, \texttt{horizon\_months}, et les 5 variables
macro projet\'ees.


\newpage
% ═══════════════════════════════════════════════════════════
\section{Orchestration : combinaison des 5 couches}
% ═══════════════════════════════════════════════════════════

\subsection{Classe CROAnalyst}

La classe \texttt{CROAnalyst} (fichier \texttt{orchestrator.py}) pilote le pipeline
complet. Son constructeur re\c{c}oit les r\'esultats credit et PE, les variables macro,
et un seuil ECL optionnel pour le RST personnalis\'e (FR54).

\begin{verbatim}
class CROAnalyst:
    def __init__(self, result_credit, result_pe, macro_params,
                 target_ecl=None)
    def analyze(self) -> AnalyticsState
\end{verbatim}

\subsection{Flux d'ex\'ecution}

L'appel \`a \texttt{analyze()} d\'eclenche le pipeline suivant :

\begin{enumerate}[label=\textbf{\arabic*.}]
    \item \textbf{Passe 1} --- \texttt{\_run\_layers(regime=None, pass=1)} :
    \begin{itemize}[nosep]
        \item Couche~1 : \texttt{analyze\_crossings} $\to$ asym\'etries + contributions
        \item Couche~2 : \texttt{decompose\_proportional} $\to$ allocation proportionnelle + facteurs
        \item Couche~3 : Construction du proxy ECL logit, puis \texttt{find\_tipping\_points}
              et \texttt{reverse\_stress\_test}
        \item Couche~4 : \texttt{classify\_regime} $\to$ r\'egime d\'etect\'e
        \item Couche~5 : \texttt{project\_trajectories}, \texttt{compute\_risk\_appetite},
              \texttt{compute\_early\_warning}
    \end{itemize}

    \item \textbf{D\'etection du r\'egime} --- extraction du r\'egime d\'etect\'e en passe~1
          (\texttt{state\_p1.regime}).

    \item \textbf{Passe 2} --- \texttt{\_run\_layers(regime=regime\_p1, pass=2)} :
    \begin{itemize}[nosep]
        \item Identique \`a la passe~1, mais le r\'egime conditionne :
        \begin{itemize}[nosep]
            \item Couche~2 : multiplicateurs r\'egime ($m_{v,r}$) sur l'attribution factorielle
            \item Couche~5 : drifts additionnels ($\delta_{v,r}$) sur les trajectoires O-U
        \end{itemize}
    \end{itemize}

    \item \textbf{Synth\`ese} --- apr\`es la passe~2 :
    \begin{itemize}[nosep]
        \item \texttt{\_generate\_recommendations} $\to$ liste de \texttt{Recommendation}
        \item \texttt{\_generate\_narrative} $\to$ synth\`ese texte structur\'ee
    \end{itemize}
\end{enumerate}

\subsection{Construction du proxy ECL dans l'orchestrateur}

Le proxy ECL est construit au sein de \texttt{\_run\_layers} \`a partir du portefeuille
r\'eel (pas de constantes hardcod\'ees) :

\begin{enumerate}[nosep]
    \item $\text{ECL}_{\text{base}} = \sum_i \text{ecl\_weighted}_i$ depuis \texttt{result\_credit}
    \item $\text{EAD}_{\text{total}} = \sum_i \text{ead}_i$
    \item $r_{\text{base}} = \text{ECL}_{\text{base}} / \text{EAD}_{\text{total}}$
    \item $\text{logit}_{\text{base}} = \text{logit}(\text{clip}(r_{\text{base}}, 10^{-6}, 0{,}99))$
    \item Mapping variable $\to$ attribut \texttt{SectorConfig} :
          \texttt{unemployment\_rate} $\to$ \texttt{unemployment\_sensitivity\_credit}, etc.
    \item Variables invers\'ees : \texttt{gdp\_growth} et \texttt{hpi\_growth}
          (baisse = stress positif)
\end{enumerate}

Le seuil de breach pour le RST utilise le capital r\'eel :
$\text{Capital} = \text{RWA}_{\text{credit}} \times \text{CET1}_{\text{target}}$
(coh\'erent avec \texttt{comparator.py}).

\subsection{G\'en\'eration des recommandations (FR31)}

La m\'ethode \texttt{\_generate\_recommendations} construit une cha\^ine de raisonnement
compl\`ete :

\begin{figure}[htbp]
\centering
\begin{tikzpicture}[
    step/.style={
        rectangle, rounded corners=4pt, minimum width=2.6cm, minimum height=0.8cm,
        draw=steelblue, thick, fill=steelblue!10, font=\scriptsize\bfseries,
        text=darkblue, align=center
    },
    arr/.style={-{Stealth[length=4pt]}, thick, steelblue},
]
\node[step] (reg) at (0,0) {R\'egime\\d\'etect\'e};
\node[step] (tri) at (3,0) {D\'eclencheur\\(risk appetite)};
\node[step] (eul) at (6,0) {Driver\\proportionnel};
\node[step] (mac) at (9,0) {Facteur\\macro};
\node[step] (rst) at (12,0) {Distance\\RST ($\sigma$)};
\node[step] (con) at (15,0) {Confiance\\(H/M/L)};

\draw[arr] (reg) -- (tri);
\draw[arr] (tri) -- (eul);
\draw[arr] (eul) -- (mac);
\draw[arr] (mac) -- (rst);
\draw[arr] (rst) -- (con);
\end{tikzpicture}
\caption{Cha\^ine de raisonnement des recommandations}
\label{fig:reco_chain}
\end{figure}

Le statut risk appetite pr\'edominant est d\'etermin\'e par le nombre de secteurs en rouge :

\begin{itemize}[nosep]
    \item $> 3$ secteurs rouge $\Rightarrow$ statut \textcolor{accentred}{\textbf{rouge}} :
          r\'eduire l'exposition + renforcer les provisions
    \item $1$--$3$ secteurs rouge $\Rightarrow$ statut \textcolor{accentorange}{\textbf{ambre}} :
          surveiller les indicateurs + plan de contingence
    \item $0$ secteur rouge $\Rightarrow$ statut \textcolor{accentgreen}{\textbf{vert}} :
          maintenir l'allocation actuelle
\end{itemize}

Chaque \texttt{Recommendation} porte :
\texttt{action}, \texttt{regime}, \texttt{trigger}, \texttt{proportional\_driver},
\texttt{macro\_factor}, \texttt{risk\_appetite\_status}, \texttt{rst\_distance},
\texttt{confidence}, et \texttt{alternatives}.

\subsection{G\'en\'eration de la synth\`ese narrative (FR32)}

La synth\`ese narrative est produite par \texttt{\_generate\_narrative} et inclut :

\begin{enumerate}[nosep]
    \item R\'egime d\'etect\'e avec sa probabilit\'e
    \item Top~3 des asym\'etries sectorielles
    \item Driver proportionnel dominant (share du risque total)
    \item Distance RST en $\sigma$ (Mahalanobis)
    \item Recommandation principale avec son niveau de confiance
\end{enumerate}

\begin{warningbox}[Performance]
Le pipeline complet (2 passes) s'ex\'ecute en $\sim 1{,}1$~s gr\^ace au proxy ECL logit.
Sans proxy, le recalcul complet d'\texttt{ECLCalculator} pour chaque point du balayage RST
(90 it\'erations $\times$ 5 variables pour les tipping points + 90 it\'erations pour le RST
joint) rendrait l'analyse interactive impossible.
\end{warningbox}


\newpage
% ═══════════════════════════════════════════════════════════
\section{Module de commentaire : LocalCROAnalyst}
% ═══════════════════════════════════════════════════════════

\subsection{R\^ole et positionnement}

Le module \texttt{analytics/ai\_analyst.py} impl\'emente la classe
\texttt{LocalCROAnalyst}, un moteur d'analyse statistique embarqu\'e (sans d\'ependance
externe, pas d'API) qui produit un rapport CRO narratif de niveau Direction des Risques.

Contrairement au pipeline 5 couches (qui produit des m\'etriques structur\'ees dans
\texttt{AnalyticsState}), le \texttt{LocalCROAnalyst} g\'en\`ere :

\begin{itemize}
    \item Un \textbf{rapport complet} en Markdown structur\'e (8 sections), via
          \texttt{generate\_full\_report()}.
    \item Un \textbf{briefing ex\'ecutif} en prose professionnelle, via
          \texttt{generate\_executive\_briefing()}.
\end{itemize}

\subsection{Architecture du rapport (8 sections)}

\begin{table}[htbp]
\centering
\caption{Sections du rapport CRO et contenu}
\begin{tabular}{clp{8cm}}
\toprule
\textbf{N} & \textbf{Section} & \textbf{Contenu} \\
\midrule
1 & Synth\`ese ex\'ecutive & Niveau de risque (score 0--10), ECL total,
    coverage, points d'attention \\
2 & Analyse ECL & D\'ecomposition waterfall, attribution par driver
    et segment \\
3 & Sensibilit\'e macro & Tableau des 5 variables, canaux de transmission,
    effets combin\'es \\
4 & Qualit\'e de cr\'edit & Distribution staging, matrice de transition,
    vuln\'erabilit\'e par segment \\
5 & Gouvernance mod\`eles & AUC/Gini/KS par mod\`ele, PSI, backtesting
    walk-forward, SHAP \\
6 & Concentration & HHI segments et produits, contribution ECL par
    segment \\
7 & Forward-looking & Risques \'emergents, pipeline Stage~2,
    sc\'enarios prospectifs \\
8 & Recommandations & Recommandations prioris\'ees (HAUTE/MOYENNE/STANDARD)
    avec KPI de suivi \\
\bottomrule
\end{tabular}
\label{tab:report_sections}
\end{table}

\subsection{Scoring du risque global}

Le scoring agrege 4 dimensions sur un total de 10 points :

\begin{formulebox}[Score de risque global (0--10)]
\begin{equation}
    S = S_{\text{ECL}} + S_{\text{Stage2}} + S_{\text{PSI}} + S_{\text{Macro}}
    \label{eq:risk_score}
\end{equation}
\begin{itemize}[nosep]
    \item $S_{\text{ECL}} \in \{0, 1, 3\}$ : 3 si $|\Delta\text{ECL}| > 15\%$, 1 si $> 5\%$
    \item $S_{\text{Stage2}} \in \{0, 1, 2\}$ : 2 si $\%\text{S2} > 20\%$, 1 si $> 15\%$
    \item $S_{\text{PSI}} \in \{0, 1, 2\}$ : 2 si PSI $> 0{,}25$, 1 si $> 0{,}10$
    \item $S_{\text{Macro}} \in \{0, 1, 2\}$ : 2 si $\geq 3$ variables adverses, 1 si $\geq 1$
\end{itemize}
Classification : $S \geq 5$ = \textcolor{accentred}{ELEVE}, $S \geq 2$ = \textcolor{accentorange}{MODERE}, sinon \textcolor{accentgreen}{BON}.
\end{formulebox}

\subsection{\'Elasticit\'es macro-cr\'edit}

Le \texttt{LocalCROAnalyst} calcule les \'elasticit\'es macro-cr\'edit par segment
et identifie les canaux de transmission dominants. Pour chaque segment $s$
et variable macro $v$ :

\begin{equation}
    \text{mult}_{\text{PD}} = 1 + \frac{\max(0, \Delta_v^{\text{adverse}})}{100}
    \cdot \gamma_{s,v} \cdot 10
    \label{eq:pd_mult}
\end{equation}

Le driver macro dominant d'un segment est la variable avec le plus grand
multiplicateur de PD.


\newpage
% ═══════════════════════════════════════════════════════════
\section{Synth\`ese math\'ematique}
% ═══════════════════════════════════════════════════════════

\subsection{R\'ecapitulatif des formules cl\'es}

\begin{longtable}{clp{8cm}}
\caption{Formules cl\'es de l'AI Analyst par couche} \\
\toprule
\textbf{Couche} & \textbf{Formule} & \textbf{R\'ef\'erence} \\
\midrule
\endfirsthead
\midrule
\textbf{Couche} & \textbf{Formule} & \textbf{R\'ef\'erence} \\
\midrule
\endhead
1 & $\text{Asym}_s = \ell_s^{\text{PE}} - \ell_s^{\text{Credit}}$ & \'Eq.\,(\ref{eq:asymmetry}) \\
1 & $c_{s,k} = R_{s,k} / R_{\text{total}}$ & \'Eq.\,(\ref{eq:marginal}) \\
1 & $\text{EW}_s = \sum_j w_j \cdot \tilde{x}_{s,j}$ & \'Eq.\,(\ref{eq:early_warning}) \\
\midrule
2 & $w_s^k = R_s^k / R_{\text{total}}$ (full allocation) & \'Eq.\,(\ref{eq:proportional}) \\
2 & $A_v^k = \sum_s (|\Delta_v| / \delta_v) \cdot \gamma_{s,v}^k \cdot R_s^k / R^k$ & \'Eq.\,(\ref{eq:attribution}) \\
2 & $A_v^{k,\text{ajust}} = A_v^k \cdot m_{v,r}$ & \'Eq.\,(\ref{eq:regime_mult}) \\
\midrule
3 & $\text{ECL}_{\text{proxy}} = \sigma(\text{logit}(r) + \Lambda \cdot \text{stress}) \cdot \text{EAD}$ & \'Eq.\,(\ref{eq:ecl_proxy}) \\
3 & $d_M = \sqrt{\mathbf{x}^\top \boldsymbol{\Sigma}^{-1} \mathbf{x}}$ (Mahalanobis) & \'Eq.\,(\ref{eq:mahalanobis}) \\
\midrule
4 & $z_r = \mathbf{s} \cdot \boldsymbol{\sigma}_r$ (produit scalaire) & \'Eq.\,(\ref{eq:dot_product}) \\
4 & $P(r) = \exp(T z_r) / \sum \exp(T z_{r'})$ (softmax) & \'Eq.\,(\ref{eq:softmax}) \\
\midrule
5 & $X_{t+1} = X_t + \kappa(\theta - X_t)\Delta t + \sigma\sqrt{\Delta t}\,\varepsilon$ (O-U) & \'Eq.\,(\ref{eq:ou_discrete}) \\
5 & $X^{\text{ajust}} = X^{\text{proj}} + \delta_{v,r} \cdot T/12$ & \'Eq.\,(\ref{eq:regime_drift}) \\
\bottomrule
\end{longtable}

\subsection{Flux de donn\'ees complet}

\begin{figure}[htbp]
\centering
\begin{tikzpicture}[
    box/.style={
        rectangle, rounded corners=3pt, minimum width=3cm, minimum height=0.7cm,
        draw=steelblue, fill=steelblue!8, font=\scriptsize, text=darkblue, align=center
    },
    arr/.style={-{Stealth[length=4pt]}, thick, gray},
]
% Inputs
\node[box, fill=accentorange!15, draw=accentorange] (in1) at (0,3) {ECLCalculator};
\node[box, fill=accentorange!15, draw=accentorange] (in2) at (0,2) {PECalculator};
\node[box, fill=accentorange!15, draw=accentorange] (in3) at (0,1) {macro\_params};

% Orchestrateur
\node[box, minimum width=4cm, fill=darkblue!15, draw=darkblue] (orc) at (5,2) {CROAnalyst\\orchestrator.py};

% Couches
\node[box] (c1) at (10,4) {Couche 1\\Croisements};
\node[box] (c2) at (10,3) {Couche 2\\Allocation};
\node[box] (c3) at (10,2) {Couche 3\\RST};
\node[box] (c4) at (10,1) {Couche 4\\R\'egime};
\node[box] (c5) at (10,0) {Couche 5\\Prospective};

% Output
\node[box, fill=accentgreen!15, draw=accentgreen!70!black, minimum width=3.5cm] (out) at (15,2) {AnalyticsState\\+ Narrative\\+ Recommendations};

% Arrows
\draw[arr] (in1) -- (orc);
\draw[arr] (in2) -- (orc);
\draw[arr] (in3) -- (orc);

\draw[arr] (orc) -- (c1);
\draw[arr] (orc) -- (c2);
\draw[arr] (orc) -- (c3);
\draw[arr] (orc) -- (c4);
\draw[arr] (orc) -- (c5);

\draw[arr] (c1) -- (out);
\draw[arr] (c2) -- (out);
\draw[arr] (c3) -- (out);
\draw[arr] (c4) -- (out);
\draw[arr] (c5) -- (out);

\end{tikzpicture}
\caption{Flux de donn\'ees du pipeline CRO}
\label{fig:dataflow}
\end{figure}


\newpage
% ═══════════════════════════════════════════════════════════
\section{Annexes}
% ═══════════════════════════════════════════════════════════

\subsection{Sc\'enarios pr\'ed\'efinis}

Les 5 sc\'enarios pr\'ed\'efinis (dropdown dans le dashboard) sont d\'efinis dans
\texttt{PREDEFINED\_SCENARIOS} (fichier \texttt{config.py}) :

\begin{table}[htbp]
\centering
\caption{Sc\'enarios pr\'ed\'efinis (\texttt{config.py})}
\begin{tabular}{lccccc}
\toprule
\textbf{Sc\'enario} & \textbf{Taux (bp)} & \textbf{Ch\^omage*} & \textbf{PIB (\%)} & \textbf{HPI (\%)} & \textbf{Infl. (\%)} \\
\midrule
Central & +50 & 0 & 1,2 & 2,0 & 2,5 \\
Crise financi\`ere & +200 & $-4$ & $-3{,}0$ & $-15{,}0$ & $-1{,}0$ \\
Stagflation & +300 & $-2$ & $-1{,}0$ & $-5{,}0$ & 6,0 \\
Rupture techno & 0 & $+3$ & 2,0 & 5,0 & 1,0 \\
Reprise & $-100$ & 0 & 3,0 & 8,0 & 2,0 \\
\bottomrule
\end{tabular}

\vspace{0.3cm}
\footnotesize
* Convention bipolaire : signe = nature ($+$ rupture techno, $-$ crise \'eco),
$|$valeur$|$ = amplitude de hausse du ch\^omage en pp vs base.
Les deux extr\^emes augmentent le taux de ch\^omage : base $+ |$slider$|$.
\label{tab:scenarios}
\end{table}

\subsection{Matrice de covariance macro}

\begin{warningbox}[Audit RJ v2 -- matrice de consensus expert]
\textbf{Correction critique} : l'ancienne impl\'ementation g\'en\'erait 60 mois d'historique
synth\'etique via des processus AR(1) \textbf{ind\'ependants}, puis estimait $\hat{\Sigma}$
via Ledoit-Wolf. Probl\`eme de circularit\'e : des AR(1) ind\'ependants produisent une
covariance quasi-diagonale --- la r\'e-estimation est inutile.

La matrice de covariance est d\'esormais construite directement \`a partir de :
\begin{itemize}[nosep]
    \item une \textbf{matrice de corr\'elation de consensus} (\textit{Expert Judgment})
          calibr\'ee sur la litt\'erature macro-financi\`ere (loi d'Okun, courbe de Phillips,
          r\`egle de Taylor, corr\'elations FMI WEO 2000--2024) ;
    \item des \textbf{volatilit\'es historiques} par variable ($\sigma_i$ en pp, sources BCE/Eurostat) ;
    \item la formule $\Sigma_{ij} = \rho_{ij} \times \sigma_i \times \sigma_j$ ;
    \item une \textbf{r\'egularisation} (r\'etractation 10\% vers la cible diagonale)
          pour garantir la d\'efinie-positivit\'e num\'erique.
\end{itemize}
Impl\'ementation : \texttt{\_build\_macro\_covariance()} dans \texttt{config.py}.
\end{warningbox}

L'ordre des variables dans $\boldsymbol{\Sigma}$ est :

\begin{enumerate}[nosep]
    \item \texttt{unemployment\_rate}
    \item \texttt{gdp\_growth}
    \item \texttt{interest\_rate}
    \item \texttt{hpi\_growth}
    \item \texttt{inflation\_rate}
\end{enumerate}

Cette matrice capture les corr\'elations historiques entre les variables macro
(par exemple : ch\^omage/PIB n\'egativement corr\'el\'es, taux/inflation positivement
corr\'el\'es), permettant \`a la distance de Mahalanobis de p\'enaliser les sc\'enarios
incoh\'erents o\`u les variables \'evoluent dans des directions contraires \`a leurs
corr\'elations historiques.

\subsection{Fichiers source}

\begin{table}[htbp]
\centering
\caption{Correspondance fichiers/fonctions}
\begin{tabular}{lll}
\toprule
\textbf{Fichier} & \textbf{Fonction principale} & \textbf{Couche} \\
\midrule
\texttt{layer1\_crossing.py} & \texttt{analyze\_crossings()} & 1 \\
\texttt{layer2\_euler.py} & \texttt{decompose\_proportional()} & 2 \\
\texttt{layer3\_rst.py} & \texttt{find\_tipping\_points()}, \texttt{reverse\_stress\_test()} & 3 \\
\texttt{layer4\_regime.py} & \texttt{classify\_regime()} & 4 \\
\texttt{layer5\_prospective.py} & \texttt{project\_trajectories()}, \texttt{compute\_risk\_appetite()}, \texttt{compute\_early\_warning()} & 5 \\
\texttt{orchestrator.py} & \texttt{CROAnalyst.analyze()} & Orchestrateur \\
\texttt{types.py} & \texttt{AnalyticsState}, \texttt{Recommendation}, \texttt{RegimeClassification} & Types \\
\texttt{analytics/ai\_analyst.py} & \texttt{LocalCROAnalyst.generate\_full\_report()} & Commentaire \\
\texttt{config.py} & \texttt{MACRO\_MEAN\_REVERSION}, \texttt{MACRO\_COVARIANCE}, etc. & Config \\
\bottomrule
\end{tabular}
\label{tab:source_files}
\end{table}

\subsection{Constantes cl\'es}

\begin{table}[htbp]
\centering
\caption{Constantes de configuration utilis\'ees par l'AI Analyst}
\begin{tabular}{llr}
\toprule
\textbf{Constante} & \textbf{Description} & \textbf{Valeur} \\
\midrule
\texttt{LOGIT\_AMPLITUDE} & Amplitude du stress logit (calibr\'ee EBA) & $\approx 4{,}94$ \\
\texttt{\_SOFTMAX\_SCALE} & Temp\'erature softmax r\'egime & 2,0 \\
\texttt{cet1\_target} & Ratio CET1 cible & 13\% \\
\texttt{rwa\_budget} & Budget RWA & 10 Md EUR \\
\texttt{ecl\_ead\_green} & Seuil vert ECL/EAD & 2,0\% \\
\texttt{ecl\_ead\_amber} & Seuil ambre ECL/EAD & 4,0\% \\
\texttt{nav\_drawdown\_green} & Seuil vert drawdown PE & 10\% \\
\texttt{nav\_drawdown\_amber} & Seuil ambre drawdown PE & 25\% \\
\bottomrule
\end{tabular}
\label{tab:constants}
\end{table}

\end{document}
